\section{Conclusion}
\label{sec:conclusion}

This project designed and simulated a 5G microstrip patch antenna and a \(2\times1\) phased patch
array. The selected design frequency was \(f_0=2.501~\mathrm{GHz}\), the substrate was
DiClad~880-IM, and the substrate thickness was \(h=1.524~\mathrm{mm}\). The analytical
transmission-line model provided the initial patch dimensions, and CST Studio Suite was then used
to tune and evaluate the final geometry.

The final tuned single patch used \(\texttt{patchL}=38.40~\mathrm{mm}\) and
\(\texttt{insetY}=14.11~\mathrm{mm}\). The PEC baseline resonance occurred at approximately
\(2.50507~\mathrm{GHz}\), with \(S_{11,\min}\approx-39.87~\mathrm{dB}\). The peak
directivity was approximately \(7.226~\mathrm{dBi}\), and the radiation and total efficiencies
were approximately \(96.3\%\) and \(96.2\%\), respectively. A supplementary fixed-mesh
finite-conductivity copper run produced \(f_r\approx2.504~\mathrm{GHz}\),
\(S_{11,\min}\approx-39.72~\mathrm{dB}\), and directivity near \(7.23~\mathrm{dBi}\), while
the radiation and total efficiencies decreased to approximately \(92.46\%\) and \(92.38\%\).
This comparison confirmed that the PEC approximation preserved the principal resonance and
radiation behavior but overestimated efficiency by omitting finite-conductor loss.

The design was then extended to a \(2\times1\) array with an element spacing of
\(d=59.935~\mathrm{mm}\). Because mutual coupling shifted the resonance relative to the isolated
single patch, the array was retuned to \(\texttt{patchL}=36.70~\mathrm{mm}\) and
\(\texttt{insetY}=13.35~\mathrm{mm}\). The final array achieved acceptable two-port matching
near the operating frequency, with \(S_{11}\approx-11.93~\mathrm{dB}\),
\(S_{22}\approx-15.20~\mathrm{dB}\), and mutual coupling of approximately
\(S_{12}\approx S_{21}\approx-15~\mathrm{dB}\). The coupling result explains why the array
required independent tuning after the isolated element was optimized.

Beam steering was demonstrated using CST amplitude-phase post-processing. The two port farfields
were combined for equal-amplitude excitations with relative phases of \(0^\circ\), \(-90^\circ\),
and \(+90^\circ\). The \(0^\circ\) case produced the broadside reference pattern, with the
principal-plane main-beam direction remaining within approximately \(0^\circ\)--\(1^\circ\).
The \(-90^\circ\) and \(+90^\circ\) cases produced simulated beam displacements of
approximately \(20^\circ\) toward the \(\phi=0^\circ\) and \(\phi=180^\circ\) half-planes,
respectively. For an element spacing of approximately \(\lambda_0/2\), the ideal array-factor
model predicts a scan-angle magnitude near \(30^\circ\). The difference is consistent with the
finite element pattern, mutual coupling, substrate loading, finite ground plane, and practical feed
geometry included in the full-wave model.

The numerical results must be interpreted within the limitations of the CST Learning Edition
environment. The principal tuning and array simulations used PEC metallization, while the
supplementary single-patch validation used finite-conductivity lossy-metal copper. An attempted
adaptive-mesh copper run exceeded the \(20{,}000\)-cell limit before convergence and was discarded;
the successful copper comparison therefore used a fixed mesh. Formal mesh-convergence and
boundary-convergence studies were not completed. Nevertheless, the observed trends were physically
consistent: reducing the isolated-patch length shifted the resonance upward, forming the array
produced a downward resonance shift accompanied by quantifiable mutual coupling, and including
copper loss reduced efficiency without materially changing resonance, matching, directivity, or
normalized pattern shape.

A higher-fidelity validation study should repeat the final single-patch and array simulations using
progressively refined local meshes and increased open-boundary spacing. The resulting resonant
frequency, \(S\)-parameters, directivity, efficiency, and beam-steering angle should then be compared
to confirm numerical convergence. A higher-capacity solver license would also permit adaptive
finite-conductivity validation of the complete two-element array.

The public project package connects the selected design inputs, analytical patch synthesis, CST
tuning, mutual-coupling analysis, numerical-modeling limitations, and phase-controlled beam
steering into a traceable antenna-engineering workflow. Selected simulation figures, analytical
scripts, parameter tables, and reproducibility notes accompany the report. Native CST project files
are not included in the ordinary source archive because of file-size and distribution considerations.
